\documentclass[12pt]{article}
\usepackage[utf8]{inputenc}
\usepackage[T1]{fontenc}
\usepackage{hyperref}
\title{Zhuang, Liuqian Language Resources for K-12 Curriculum Development}
\author{Compiled by Claude Sonnet 5.5 (claude-sonnet-5-5) for Andr\'as Kornai}
\date{}
\begin{document}
\maketitle

\section{Introduction}
Liuqian Zhuang (ISO 639-3: zlq; Glottocode liuq1235) is a Tai-Kadai language of the Zhuang macrolanguage, spoken in China. Speaker estimates in the search results are about 406,000 to 412,000. It is reported as taught in no schools. Searches found no primers, textbooks or dictionaries specific to Liuqian Zhuang. Joshua Project lists its written status as unknown.

\section{Primary Educational Resources}
No primers, textbooks or school programmes specific to Liuqian Zhuang were found in the search results.

\section{Linguistic and Reference Materials}
\begin{itemize}
\item Overview of the Zhuang languages and their writing systems: \url{https://en.wikipedia.org/wiki/Zhuang_languages}
\end{itemize}

\section{Religious and Community Resources}
\begin{itemize}
\item Joshua Project people-group profile. It says a Bible translation is needed: \url{https://joshuaproject.net/people_groups/19802}
\item Joshua Project language page: \url{https://www.joshuaproject.net/index.php/languages/zlq}
\end{itemize}

\section{Digital and Audio Resources}
\begin{itemize}
\item Global Recordings Network page for zlq, which lists no recordings: \url{https://globalrecordings.net/en/language/zlq}
\end{itemize}

\section{Curriculum Development Context}
With roughly 400,000 speakers, the community is fairly large, but no written or school materials were found. A curriculum would start with documentation and orthography work, and could draw on the wider Zhuang literacy tradition. How closely that tradition matches Liuqian was not established.

\end{document}